\section{Sites}
\label{section-sites-definitions}

\noindent
Our notion of a site uses the following type of structures.

\begin{definition}
\label{definition-family-morphisms-fixed-target}
Let $\mathcal{C}$ be a category, see
Conventions, Section \ref{conventions-section-categories}.
A {\it family of morphisms with fixed target} in $\mathcal{C}$ is
given by an object $U \in \Ob(\mathcal{C})$, a set $I$ and
for each $i\in I$ a morphism $U_i \to U$ of $\mathcal{C}$ with target $U$.
We use the notation $\{U_i \to U\}_{i\in I}$ to indicate this.
\end{definition}

\noindent
It can happen that the set $I$ is empty! This
notation is meant to suggest an open covering as in topology.

\begin{definition}
\label{definition-site}
A {\it site}\footnote{This notation differs from that of \cite{SGA4}, as
explained in the introduction.} is given by a category $\mathcal{C}$ and a set
$\text{Cov}(\mathcal{C})$ of families of morphisms with fixed target
$\{U_i \to U\}_{i \in I}$, called {\it coverings of $\mathcal{C}$},
satisfying the following axioms
\begin{enumerate}
\item If $V \to U$ is an isomorphism then $\{V \to U\} \in
\text{Cov}(\mathcal{C})$.
\item If $\{U_i \to U\}_{i\in I} \in \text{Cov}(\mathcal{C})$ and for each
$i$ we have $\{V_{ij} \to U_i\}_{j\in J_i} \in \text{Cov}(\mathcal{C})$, then
$\{V_{ij} \to U\}_{i \in I, j\in J_i} \in \text{Cov}(\mathcal{C})$.
\item If $\{U_i \to U\}_{i\in I}\in \text{Cov}(\mathcal{C})$
and $V \to U$ is a morphism of $\mathcal{C}$ then $U_i \times_U V$
exists for all $i$ and
$\{U_i \times_U V \to V \}_{i\in I} \in \text{Cov}(\mathcal{C})$.
\end{enumerate}
\end{definition}

\noindent
Clarifications. In axiom (1) we require there should be a covering
$\{U_i \to U\}_{i \in I}$ of $\mathcal{C}$ such that $I = \{i\}$
is a singleton set and such that the morphism $U_i \to U$ is equal to
the morphism $V \to U$ given in (1). In the following we often denote
$\{V \to U\}$ a family of morphisms with fixed target whose index set
is a singleton. In axiom (3) we require the existence of the covering
for some choice of the fibre products $U_i \times_U V$ for $i \in I$.

\begin{remark}
\label{remark-no-big-sites}
(On set theoretic issues -- skip on a first reading.)
The main reason for introducing sites is to study the
category of sheaves on a site, because it is the generalization
of the category of sheaves on a topological space that has
been so important in algebraic geometry. In order to avoid thinking
about things like ``classes of classes'' and so on, we will
not allow sites to be ``big'' categories, in contrast to what
we do for categories and $2$-categories.

\medskip\noindent
Suppose that $\mathcal{C}$ is a category and that
$\text{Cov}(\mathcal{C})$ is a proper class of coverings
satisfying (1), (2) and (3) above. We will not allow this as a
site either, mainly because we are going to take limits over coverings.
However, there are several natural
ways to replace $\text{Cov}(\mathcal{C})$ by a set of coverings
or a slightly different structure
that give rise to the same category of sheaves. For example:
\begin{enumerate}
\item In Sets, Section \ref{sets-section-coverings-site}
we show how to pick a suitable set of
coverings that gives the same category of sheaves.
\item Another thing we can do is to take the associated topology
(see Definition \ref{definition-topology-associated-site}).
The resulting topology on $\mathcal{C}$ has the same category of sheaves.
Two topologies have the same categories of sheaves if and only if
they are equal, see Theorem \ref{theorem-topology-and-topos}.
A topology on a category is given by a choice of sieves on objects.
The collection of all possible sieves and even all possible
topologies on $\mathcal{C}$ is a set.
\item We could also slightly modify the notion of a site, see
Remark \ref{remark-shrink-coverings} below, and end up with a
canonical set of coverings.
\end{enumerate}
Each of these solutions has some minor drawback. For the first, one has
to check that constructions later on do not depend on the choice
of the set of coverings. For the second, one has to learn about topologies
and redo many of the arguments for sites. For the third, see
the last sentence of Remark \ref{remark-shrink-coverings}.

\medskip\noindent
Our approach will be to work with sites as in Definition \ref{definition-site}
above. Given a category $\mathcal{C}$ with a proper class of coverings
as above, we will replace this by a set of coverings producing a site using
Sets, Lemma \ref{sets-lemma-coverings-site}. It is shown in
Lemma \ref{lemma-choice-set-coverings-immaterial} below that the resulting
category of sheaves (the topos) is independent of this choice. We leave it to
the reader to use one of the other two strategies to deal with these issues if
he/she so desires.
\end{remark}

\begin{example}
\label{example-site-topological}
Let $X$ be a topological space. Let $X_{Zar}$ be the category whose
objects consist of all the open sets $U$ in $X$ and whose morphisms
are just the inclusion maps. That is, there is at most one morphism
between any two objects in $X_{Zar}$. Now define
$\{U_i \to U\}_{i \in I}\in \text{Cov}(X_{Zar})$ if
and only if $\bigcup U_i = U$.
Conditions (1) and (2) above are clear, and (3) is also
clear once we realize that in $X_{Zar}$ we have
$U \times V = U \cap V$. Note that in particular the empty
set has to be an element of $X_{Zar}$ since otherwise
this would not work in general. Furthermore, it is equally important,
as we will see later, to allow the {\it empty covering of the empty
set as a covering}!
We turn $X_{Zar}$ into a site
by choosing a suitable set of coverings
$\text{Cov}(X_{Zar})_{\kappa, \alpha}$ as in
Sets, Lemma \ref{sets-lemma-coverings-site}.
Presheaves and sheaves (as defined below)
on the site $X_{Zar}$ agree exactly with the usual notion of
a presheaves and sheaves on a topological space, as defined
in Sheaves, Section \ref{sheaves-section-introduction}.
\end{example}

\begin{example}
\label{example-site-on-group}
Let $G$ be a group. Consider the category $G\textit{-Sets}$
whose objects are sets $X$ with a left $G$-action, with
$G$-equivariant maps as the morphisms. An important example
is ${}_GG$ which is the $G$-set whose underlying set is $G$ and
action given by left multiplication. This category has
fiber products, see Categories, Section
\ref{categories-section-example-fibre-products}.
We declare $\{\varphi_i : U_i \to U\}_{i\in I}$ to be
a covering if $\bigcup_{i\in I} \varphi_i(U_i) = U$.
This gives a class of coverings on $G\textit{-Sets}$
which is easily seen to satisfy conditions (1), (2), and (3)
of Definition \ref{definition-site}. The result is not a
site since both the collection of objects of the underlying category and
the collection of coverings form a proper class.
We first replace by $G\textit{-Sets}$ by a
full subcategory $G\textit{-Sets}_\alpha$ as in Sets,
Lemma \ref{sets-lemma-sets-with-group-action}.
After this the site
$(G\textit{-Sets}_\alpha,
\text{Cov}_{\kappa, \alpha'}(G\textit{-Sets}_\alpha))$
gotten by suitably restricting the collection of coverings
as in Sets, Lemma \ref{sets-lemma-coverings-site} will be
denoted $\mathcal{T}_G$.

\medskip\noindent
As a special case, if the group $G$ is countable, then we can let
$\mathcal{T}_G$ be the category of countable $G$-sets and coverings
those jointly surjective families of morphisms
$\{\varphi_i : U_i \to U\}_{i \in I}$ such that $I$ is countable.
\end{example}

\begin{example}
\label{example-indiscrete}
Let $\mathcal{C}$ be a category. There is a canonical way to turn this
into a site where $\{f : V \to U \mid f\text{ is an isomorphism}\}$
are the coverings of $U$.
Sheaves on this site are the presheaves on $\mathcal{C}$.
This corresponding topology is called the {\it chaotic} or
{\it indiscrete topology}.
\end{example}















